\chapter{Complete System Call Reference}
\label{app:syscalls}

This appendix provides an exhaustive, classified reference of the Linux System Calls implemented by Hydra's binary compatibility layer (\texttt{src/compat/linux/}). Calls adhere to the System V AMD64 ABI: the system call number is passed in \texttt{RAX}, arguments in \texttt{RDI}, \texttt{RSI}, \texttt{RDX}, \texttt{R10}, \texttt{R8}, and \texttt{R9}, and return values in \texttt{RAX} (negative values denote \texttt{-errno}).

% ==============================================================================
\section{File and Directory Operations}
\label{sec:ref_file_ops}
% ==============================================================================

\begin{table}[htbp]
\centering
\small
\begin{tabularx}{\textwidth}{rllX}
\toprule
\textbf{Syscall} & \textbf{Name} & \textbf{Arguments (RDI..R9)} & \textbf{Semantics and Description} \\
\midrule
0 & \texttt{read} & \texttt{fd, buf, count} & Reads up to \texttt{count} bytes from \texttt{fd} into user buffer \texttt{buf}. \\
1 & \texttt{write} & \texttt{fd, buf, count} & Writes up to \texttt{count} bytes from \texttt{buf} to \texttt{fd}. \\
2 & \texttt{open} & \texttt{path, flags, mode} & Opens file specified by \texttt{path} with flags (\texttt{O\_RDONLY}, etc.). \\
3 & \texttt{close} & \texttt{fd} & Closes an active file descriptor handle. \\
4 & \texttt{stat} & \texttt{path, statbuf} & Retrieves file metadata by pathname. \\
5 & \texttt{fstat} & \texttt{fd, statbuf} & Retrieves file metadata by open descriptor handle. \\
6 & \texttt{lstat} & \texttt{path, statbuf} & Retrieves metadata without following symlinks. \\
8 & \texttt{lseek} & \texttt{fd, offset, whence} & Repositions file offset (\texttt{SEEK\_SET}, \texttt{SEEK\_CUR}, \texttt{SEEK\_END}). \\
16 & \texttt{ioctl} & \texttt{fd, cmd, arg} & Device-specific control (DRM graphics, terminal termios). \\
17 & \texttt{pread64} & \texttt{fd, buf, count, off} & Reads from specified offset without modifying file offset. \\
18 & \texttt{pwrite64} & \texttt{fd, buf, count, off} & Writes to specified offset without modifying file offset. \\
19 & \texttt{readv} & \texttt{fd, iov, iovcnt} & Scatter-read into array of buffer descriptors. \\
20 & \texttt{writev} & \texttt{fd, iov, iovcnt} & Gather-write from array of buffer descriptors. \\
21 & \texttt{access} & \texttt{path, mode} & Checks caller's real permissions for a file path. \\
22 & \texttt{pipe} & \texttt{pipefd[2]} & Creates a unidirectional FIFO interprocess data channel. \\
32 & \texttt{dup} & \texttt{oldfd} & Duplicates file descriptor to the lowest available integer. \\
33 & \texttt{dup2} & \texttt{oldfd, newfd} & Duplicates file descriptor to a specific integer handle. \\
72 & \texttt{fcntl} & \texttt{fd, cmd, arg} & Manipulates descriptor flags (\texttt{F\_SETFL}, \texttt{O\_NONBLOCK}). \\
78 & \texttt{getdents} & \texttt{fd, dirent, count} & Reads directory entries into user buffer. \\
79 & \texttt{getcwd} & \texttt{buf, size} & Retrieves absolute working directory path string. \\
80 & \texttt{chdir} & \texttt{path} & Changes current process working directory. \\
83 & \texttt{mkdir} & \texttt{path, mode} & Creates a new directory node on the ZFS dataset. \\
84 & \texttt{rmdir} & \texttt{path} & Deletes an empty directory node. \\
87 & \texttt{unlink} & \texttt{path} & Unlinks (deletes) a named file node. \\
88 & \texttt{symlink} & \texttt{target, linkpath} & Creates a symbolic link referencing \texttt{target}. \\
89 & \texttt{readlink} & \texttt{path, buf, bufsiz} & Reads the target path contained inside a symlink. \\
257 & \texttt{openat} & \texttt{dfd, path, flags, mode} & Opens file relative to directory descriptor \texttt{dfd}. \\
258 & \texttt{mkdirat} & \texttt{dfd, path, mode} & Creates directory relative to directory descriptor. \\
263 & \texttt{unlinkat} & \texttt{dfd, path, flags} & Unlinks file relative to directory descriptor. \\
267 & \texttt{readlinkat}& \texttt{dfd, path, buf, size} & Reads symlink relative to directory descriptor. \\
268 & \texttt{fchmodat} & \texttt{dfd, path, mode, flags} & Modifies file permissions relative to directory descriptor. \\
269 & \texttt{faccessat}& \texttt{dfd, path, mode, flags} & Checks permissions relative to directory descriptor. \\
293 & \texttt{pipe2} & \texttt{pipefd[2], flags} & Creates pipe with \texttt{O\_CLOEXEC} and \texttt{O\_NONBLOCK} flags. \\
332 & \texttt{statx} & \texttt{dfd, path, flags, mask, buf} & Extended file status reporting with 64-bit timestamps. \\
\bottomrule
\end{tabularx}
\caption{Linux File and Directory System Calls Supported in Hydra.}
\label{tab:syscall_file}
\end{table}

% ==============================================================================
\section{Memory Management System Calls}
\label{sec:ref_mem_ops}
% ==============================================================================

\begin{table}[htbp]
\centering
\small
\begin{tabularx}{\textwidth}{rllX}
\toprule
\textbf{Syscall} & \textbf{Name} & \textbf{Arguments (RDI..R9)} & \textbf{Semantics and Description} \\
\midrule
9 & \texttt{mmap} & \texttt{addr, len, prot, flags, fd, off} & Maps physical frames, files, or anonymous RAM into VM. \\
10 & \texttt{mprotect} & \texttt{addr, len, prot} & Changes memory protection flags (\texttt{PROT\_READ/WRITE/EXEC}). \\
11 & \texttt{munmap} & \texttt{addr, len} & Unmaps virtual memory pages from process address space. \\
12 & \texttt{brk} & \texttt{addr} & Adjusts process heap break pointer. \\
25 & \texttt{mremap} & \texttt{old\_addr, old\_len, new\_len, flags} & Expands, shrinks, or moves existing virtual memory mapping. \\
26 & \texttt{msync} & \texttt{addr, len, flags} & Synchronizes dirty mapped pages back to storage media. \\
27 & \texttt{mincore} & \texttt{addr, len, vec} & Determines whether pages reside in physical RAM. \\
28 & \texttt{madvise} & \texttt{addr, len, advice} & Provides memory access pattern hints to the page allocator. \\
319 & \texttt{memfd\_create} & \texttt{name, flags} & Creates an anonymous in-memory file for zero-copy IPC. \\
\bottomrule
\end{tabularx}
\caption{Linux Memory Management System Calls Supported in Hydra.}
\label{tab:syscall_mem}
\end{table}

% ==============================================================================
\section{Process Lifecycle and Credential System Calls}
\label{sec:ref_proc_ops}
% ==============================================================================

\begin{table}[htbp]
\centering
\small
\begin{tabularx}{\textwidth}{rllX}
\toprule
\textbf{Syscall} & \textbf{Name} & \textbf{Arguments (RDI..R9)} & \textbf{Semantics and Description} \\
\midrule
39 & \texttt{getpid} & None & Returns current process ID (\texttt{PID}). \\
56 & \texttt{clone} & \texttt{flags, sp, ptid, ctid, tls} & Spawns thread or child process with shared resources. \\
57 & \texttt{fork} & None & Clones process using Copy-on-Write page tables. \\
58 & \texttt{vfork} & None & Spawns child, suspending parent until child calls \texttt{exec}. \\
59 & \texttt{execve} & \texttt{file, argv, envp} & Replaces process image with newly loaded ELF binary. \\
60 & \texttt{exit} & \texttt{status} & Terminates calling thread with exit code. \\
61 & \texttt{wait4} & \texttt{pid, wstatus, options, rusage} & Suspends until specified child changes state or exits. \\
102 & \texttt{getuid} & None & Returns real user ID (\texttt{UID}). \\
104 & \texttt{getgid} & None & Returns real group ID (\texttt{GID}). \\
105 & \texttt{setuid} & \texttt{uid} & Sets real, effective, and saved user ID. \\
106 & \texttt{setgid} & \texttt{gid} & Sets real, effective, and saved group ID. \\
107 & \texttt{geteuid} & None & Returns effective user ID (\texttt{EUID}). \\
108 & \texttt{getegid} & None & Returns effective group ID (\texttt{EGID}). \\
109 & \texttt{setpgid} & \texttt{pid, pgid} & Sets process group ID for terminal job control. \\
110 & \texttt{getppid} & None & Returns parent process ID. \\
111 & \texttt{getpgrp} & None & Returns process group ID of calling process. \\
112 & \texttt{setsid} & None & Creates a new session and detaches controlling terminal. \\
158 & \texttt{arch\_prctl} & \texttt{code, addr} & Sets processor architectural state (sets \texttt{FS\_BASE} TLS). \\
186 & \texttt{gettid} & None & Returns thread ID (\texttt{TID}). \\
231 & \texttt{exit\_group}& \texttt{status} & Terminates all threads in current thread group. \\
\bottomrule
\end{tabularx}
\caption{Linux Process and Credential System Calls Supported in Hydra.}
\label{tab:syscall_proc}
\end{table}

% ==============================================================================
\section{Networking and Socket System Calls}
\label{sec:ref_net_ops}
% ==============================================================================

\begin{table}[htbp]
\centering
\small
\begin{tabularx}{\textwidth}{rllX}
\toprule
\textbf{Syscall} & \textbf{Name} & \textbf{Arguments (RDI..R9)} & \textbf{Semantics and Description} \\
\midrule
41 & \texttt{socket} & \texttt{domain, type, protocol} & Creates communication socket (\texttt{AF\_INET}, \texttt{AF\_UNIX}). \\
42 & \texttt{connect} & \texttt{fd, addr, addrlen} & Initiates connection to destination IP/port or UNIX socket. \\
43 & \texttt{accept} & \texttt{fd, addr, addrlen} & Accepts incoming connection on listening socket. \\
44 & \texttt{sendto} & \texttt{fd, buf, len, flags, dest, dlen} & Transmits datagram to destination endpoint. \\
45 & \texttt{recvfrom} & \texttt{fd, buf, len, flags, src, slen} & Receives datagram and captures source address. \\
46 & \texttt{sendmsg} & \texttt{fd, msg, flags} & Transmits message with scatter-gather buffers and ancillary data. \\
47 & \texttt{recvmsg} & \texttt{fd, msg, flags} & Receives message with scatter-gather buffers and ancillary data. \\
48 & \texttt{shutdown} & \texttt{fd, how} & Shuts down reading, writing, or both on a duplex socket. \\
49 & \texttt{bind} & \texttt{fd, addr, addrlen} & Binds local IP address and port to a socket. \\
50 & \texttt{listen} & \texttt{fd, backlog} & Marks socket as passive listening server. \\
51 & \texttt{getsockname}& \texttt{fd, addr, addrlen} & Retrieves bound local address of socket. \\
52 & \texttt{getpeername}& \texttt{fd, addr, addrlen} & Retrieves remote peer address of connected socket. \\
53 & \texttt{socketpair} & \texttt{domain, type, proto, sv[2]} & Creates connected pair of anonymous UNIX domain sockets. \\
54 & \texttt{setsockopt} & \texttt{fd, level, optname, optval, optlen} & Modifies socket parameters (\texttt{SO\_REUSEADDR}, etc.). \\
55 & \texttt{getsockopt} & \texttt{fd, level, optname, optval, optlen} & Queries socket parameters. \\
288 & \texttt{accept4} & \texttt{fd, addr, addrlen, flags} & Accepts connection, setting \texttt{SOCK\_NONBLOCK/CLOEXEC}. \\
\bottomrule
\end{tabularx}
\caption{Linux Networking and Socket System Calls Supported in Hydra.}
\label{tab:syscall_net}
\end{table}

% ==============================================================================
\section{Event Multiplexing, Signals, and Timing System Calls}
\label{sec:ref_event_ops}
% ==============================================================================

\begin{table}[htbp]
\centering
\small
\begin{tabularx}{\textwidth}{rllX}
\toprule
\textbf{Syscall} & \textbf{Name} & \textbf{Arguments (RDI..R9)} & \textbf{Semantics and Description} \\
\midrule
7 & \texttt{poll} & \texttt{fds, nfds, timeout} & Waits for events on file descriptor array. \\
13 & \texttt{rt\_sigaction} & \texttt{signum, act, oldact, sigsetsize} & Configures action associated with signal. \\
14 & \texttt{rt\_sigprocmask} & \texttt{how, set, oldset, sigsetsize} & Manipulates calling thread's blocked signal mask. \\
15 & \texttt{rt\_sigreturn} & None & Returns from signal handler and restores user stack context. \\
23 & \texttt{select} & \texttt{nfds, readfds, writefds, exceptfds, timeout} & Multiplexes file descriptor bitsets. \\
24 & \texttt{sched\_yield} & None & Voluntarily surrenders remaining CPU timeslice. \\
35 & \texttt{nanosleep} & \texttt{req, rem} & Suspends thread execution for high-resolution interval. \\
62 & \texttt{kill} & \texttt{pid, sig} & Dispatches signal to specified process or group. \\
127 & \texttt{rt\_sigpending} & \texttt{set, sigsetsize} & Queries signals pending delivery to caller. \\
128 & \texttt{rt\_sigtimedwait} & \texttt{set, info, timeout, sigsetsize} & Synchronously awaits signal with timeout. \\
130 & \texttt{rt\_sigsuspend} & \texttt{mask, sigsetsize} & Temporarily replaces signal mask and suspends thread. \\
202 & \texttt{futex} & \texttt{uaddr, op, val, timeout, uaddr2, val3} & Fast user-space mutual exclusion and wait queues. \\
213 & \texttt{epoll\_create} & \texttt{size} & Allocates epoll instance descriptor. \\
228 & \texttt{clock\_gettime} & \texttt{clockid, tp} & Retrieves high-resolution time (\texttt{CLOCK\_MONOTONIC}). \\
229 & \texttt{clock\_getres} & \texttt{clockid, res} & Queries resolution of target hardware clock. \\
230 & \texttt{clock\_nanosleep} & \texttt{clockid, flags, req, rem} & High-precision sleep relative to target clock. \\
232 & \texttt{epoll\_wait} & \texttt{epfd, events, maxevents, timeout} & Waits for events on epoll set in $O(1)$ time. \\
233 & \texttt{epoll\_ctl} & \texttt{epfd, op, fd, event} & Adds, modifies, or removes descriptors from epoll set. \\
284 & \texttt{eventfd} & \texttt{initval} & Allocates interprocess event counter notification handle. \\
286 & \texttt{timerfd\_settime} & \texttt{fd, flags, new\_value, old\_value} & Arms or disarms timer file descriptor. \\
287 & \texttt{timerfd\_gettime} & \texttt{fd, curr\_value} & Queries current status of timer file descriptor. \\
290 & \texttt{eventfd2} & \texttt{initval, flags} & Allocates eventfd with \texttt{EFD\_NONBLOCK} and \texttt{EFD\_CLOEXEC}. \\
291 & \texttt{epoll\_create1} & \texttt{flags} & Allocates epoll instance with \texttt{EPOLL\_CLOEXEC}. \\
\bottomrule
\end{tabularx}
\caption{Linux Event Multiplexing, Signal, and Timing System Calls Supported in Hydra.}
\label{tab:syscall_event}
\end{table}
